\section{Refinamiento conservativo de la unidad y memoria de acarreos}
\label{lib03:refinamiento}

La completación dimensional y el registro de las operaciones admiten una
composición aritmética explícita. La arquitectura APP--TRIT--TPK proporciona
previamente el estado enriquecido, las lecturas orientadas y las
completaciones de la unidad. Sobre esas salidas se construye aquí un
transporte afín que conserva la suma normalizada y permite recuperar, por
divisiones euclídeas, toda palabra finita de acarreos canónicos. Su límite
posee cilindros de diámetro controlado; las marcas de frontera distinguen
las historias que una evaluación escalar identifica.

El bloque central, sus concatenaciones y la partición
\(1000=729+243+27+1\) son antecedentes de esta construcción. La composición
de esos objetos con el transporte afín, la inversa por residuos y su
prolongación constituye la formalización desarrollada en esta sección.
La aplicación actúa sobre un dominio aritmético declarado. Cada realización
sobre una historia enriquecida particular conserva, además, la regla
prospectiva que produce sus acarreos.

\subsection{De la norma cuadrática a la completación cúbica}
\label{lib03:completacion}

Sean \(b\geq1\) entero y \(B=b+1\geq2\). Definimos
\begin{equation}
 p_b=b^2-b+1,\qquad q_b=B^2-p_b=3b.
 \label{lib03:pb}
\end{equation}

\begin{proposition}[Transferencia de la completación]
\label{lib03:prop-completacion}
Las identidades
\begin{equation}
 Bp_b=b^3+1,\qquad Bq_b+1=B^3-b^3
 \label{lib03:cuadratica-cubica}
\end{equation}
determinan la transformación
\begin{equation}
 (p_b,q_b)\longmapsto(Bp_b-1,Bq_b+1)
       =(b^3,B^3-b^3).
 \label{lib03:transferencia-cubica}
\end{equation}
La suma pasa de \(B^2\) a \(B^3\), mientras que la suma de las
coordenadas divididas por sus respectivas escalas permanece igual a uno.
\end{proposition}

\begin{proof}
La factorización \(b^3+1=(b+1)(b^2-b+1)\) proporciona la primera
identidad. La segunda resulta de
\((b+1)^3-b^3=3b^2+3b+1\). Sumarlas prueba el cambio de escala y,
después de dividir por \(B^3\), la conservación de la unidad normalizada.
\end{proof}

Para \(b=9\),
\begin{equation}
 (73,27)\longmapsto(10\cdot73-1,10\cdot27+1)=(729,271).
 \label{lib03:caso-729}
\end{equation}
El complemento conserva la descomposición \(271=243+27+1\). En la
completación precedente, la última contribución corresponde al ápice
añadido. La corrección \((-1,+1)\) tiene suma cero, y la multiplicación
por diez transforma la escala de cien en la de mil. Para la colección
completa, el transporte de las proporciones es
\begin{equation}
 \left(\frac{p_b}{B^2},\frac{q_b}{B^2}\right)
 \longmapsto
 \left(\frac{p_b}{B^2}-\frac1{B^3},
       \frac{q_b}{B^2}+\frac1{B^3}\right).
 \label{lib03:fracciones-completacion}
\end{equation}
La cantidad conservada en esta identidad es una suma lineal. Las formas
cuadráticas y las magnitudes físicas poseen sus lectores propios, que se
compondrán posteriormente con este transporte.

\subsection{Concatenaciones del bloque central y restitución unitaria}
\label{lib03:centro}

Para \(b\geq2\), las cifras \(b-2\) y \(2\) pertenecen al alfabeto de
base \(B=b+1\). Consideremos la colección algebraica
\begin{equation}
 B_c^{(b)}=\begin{pmatrix}b-2&2\\2&b-2\end{pmatrix},\qquad
 C_B(u,v)=Bu+v.
 \label{lib03:bloque-central}
\end{equation}
El caso \(b=9\) recupera el bloque central documentado
\(\bigl(\begin{smallmatrix}7&2\\2&7\end{smallmatrix}\bigr)\), con
autovalores \(9\) y \(5\). Sus concatenaciones satisfacen
\(72+27=77+22=99\). Concatenación, espectro y posición central son
lecturas distintas del mismo antecedente y mantienen sus dominios.

La aplicación del lector a las filas de \eqref{lib03:bloque-central} da
\begin{align}
 C_B(b-2,2)&=B(b-2)+2=b(b-1)=p_b-1,\\
 C_B(2,b-2)&=2B+(b-2)=3b=q_b.
 \label{lib03:concatenaciones}
\end{align}
La suma es \(B^2-1\). La restitución de una unidad en la primera
coordenada y la transferencia siguiente producen
\begin{equation}
 (p_b-1,3b)\xrightarrow{\ +(1,0)\ }(p_b,3b)
 \xrightarrow{\ L_1\ }(b^3,B^3-b^3).
 \label{lib03:cadena-restauracion}
\end{equation}
En particular,
\begin{equation}
 (72,27)\longmapsto(73,27)\longmapsto(729,271).
 \label{lib03:cadena-central}
\end{equation}
La primera flecha incrementa la suma en uno. La segunda cambia la escala
y redistribuye una unidad con corrección de suma cero. Para representar
explícitamente la contribución previa a su incorporación, el triple
\((p_b-1,3b,1)\) ya tiene suma \(B^2\), y la aplicación
\((x,y,r)\mapsto(x+r,y,0)\) conserva esa suma. Su inversión requiere
conservar el dato \(r=1\): sobre triples arbitrarios la aplicación no
es inyectiva.

Así queda registrada la restitución sin atribuirle una conservación
cerrada de la pareja antes de incorporar la unidad. La colección en \(b\)
extiende algebraicamente el caso central recibido; la selección de los
estados del atlas continúa siendo la establecida por APP--TRIT--TPK.
La composición \eqref{lib03:cadena-restauracion} determina un acarreo
igual a uno en esa transición, conservando abierta la especificación
prospectiva de las operaciones siguientes.

\subsection{Transporte afín y dominio de positividad}
\label{lib03:afin}

\begin{definition}[Transporte con acarreo]
\label{lib03:def-transporte}
Para una base entera \(B\geq2\) y \(c\in\mathbb Z\), se define
\begin{equation}
 L_c(x,y)=(Bx-c,By+c).
 \label{lib03:lc}
\end{equation}
El entero \(c\) especifica la transferencia orientada entre las dos
coordenadas. Su signo positivo reduce la primera respecto de \(Bx\)
y aumenta la segunda en la misma cantidad.
\end{definition}

\begin{proposition}[Inversa y positividad]
\label{lib03:prop-inversa}
En la realización afín posterior \(\mathbb R^2\), con \(B,c\) fijados,
\begin{equation}
 L_c^{-1}(X,Y)=\left(\frac{X+c}{B},\frac{Y-c}{B}\right),\qquad
 X+Y=B(x+y).
 \label{lib03:inversa-afin}
\end{equation}
La preimagen es entera exactamente cuando \(X+c\) e \(Y-c\) son
divisibles por \(B\). Para \(x,y\geq0\), la imagen es no negativa
exactamente cuando
\begin{equation}
 -By\leq c\leq Bx.
 \label{lib03:positividad}
\end{equation}
\end{proposition}

\begin{proof}
La sustitución directa demuestra las dos identidades de
\eqref{lib03:inversa-afin} y su criterio de integralidad. Las
desigualdades \(Bx-c\geq0\), \(By+c\geq0\) equivalen a
\eqref{lib03:positividad}.
\end{proof}

Para el alfabeto canónico \(c\in\{0,\ldots,B-1\}\), todas las
prolongaciones son admisibles si \(x\geq1\), \(y\geq0\). En la
frontera \(x=0\), la no negatividad admite únicamente \(c=0\).
El dominio positivo y el dominio entero general quedan así
diferenciados por una condición explícita.

\Needspace{12\baselineskip}
\subsection{Recuperación canónica mediante el residuo}
\label{lib03:residuo}

\begin{theorem}[Inversa canónica por residuo]
\label{lib03:teo-residuo}
Sean \(M\geq1\) entero y \(X,Y\geq0\) enteros tales que
\(X+Y=BM\). Existe una única terna admisible \((x,y,c)\) que satisface
\[
 x+y=M,\qquad 0\leq c\leq B-1,\qquad L_c(x,y)=(X,Y).
\]
Se obtiene mediante
\begin{equation}
 c=Y\bmod B=(-X)\bmod B,\qquad
 y=\frac{Y-c}{B},\qquad x=\frac{X+c}{B}.
 \label{lib03:residuo-inversor}
\end{equation}
\end{theorem}

\begin{proof}
Toda preimagen satisface \(Y\equiv c\pmod B\). El representante en
\(\{0,\ldots,B-1\}\) es único. La congruencia \(X+Y\equiv0\pmod B\)
implica la divisibilidad de \(X+c\). La división euclídea de \(Y\)
produce \(y\geq0\), y \(X+c\geq0\) produce \(x\geq0\). Si
\(c>0\), el caso \(x=0\) exigiría \(X=-c<0\), de modo que la
terna obtenida pertenece al dominio admisible. Su suma es \(M\), y
la sustitución en \eqref{lib03:lc} recupera \((X,Y)\). La unicidad
del residuo y de los dos cocientes prueba la unicidad de la terna.
\end{proof}

La pareja entera y su escala contienen el último acarreo canónico. La
recuperación deriva de su residuo y no requiere conservar otra copia
independiente de ese acarreo. Hay \(M+1\) ternas con \(c=0\) y
\(M\) ternas en cada una de las \(B-1\) clases restantes: el total
es \(BM+1\), igual al número de parejas no negativas de suma \(BM\).

Para un acarreo entero general, escribamos
\(c=r+B\ell\), con \(0\leq r\leq B-1\). El residuo determina
\(r\), mientras que la reconstrucción completa utiliza también la
elevación \(\ell\):
\begin{equation}
 y=\frac{Y-r}{B}-\ell,\qquad
 x=\frac{X+r}{B}+\ell.
 \label{lib03:elevacion}
\end{equation}
La conservación de la elevación distingue el acarreo firmado de su
reducción residual.

\subsection{Iteración variable y memoria finita completa}
\label{lib03:memoria}

Fijemos \((x_0,y_0)\in\mathbb Z^2\), de suma \(M_0>0\), y una
secuencia de operaciones \(c_1,c_2,\ldots\). Definimos
\begin{equation}
 (x_n,y_n)=L_{c_n}\cdots L_{c_1}(x_0,y_0),\qquad
 C_n=\sum_{j=1}^n c_jB^{n-j},\qquad C_0=0.
 \label{lib03:registro-acumulado}
\end{equation}
Los acarreos se reciben de la regla de operación correspondiente; el
valor límite no interviene en su selección.

\begin{theorem}[Registro acumulado exacto]
\label{lib03:teo-acumulacion}
Para todo \(n\geq0\),
\begin{equation}
 \begin{aligned}
 x_n&=B^nx_0-C_n,& y_n&=B^ny_0+C_n,\\
 x_n+y_n&=B^nM_0,& C_{n+1}&=BC_n+c_{n+1}.
 \end{aligned}
 \label{lib03:iteracion-exacta}
\end{equation}
\end{theorem}

\begin{proof}
Para \(n=0\) las identidades son inmediatas. Si son válidas en
\(n\), la aplicación de \(L_{c_{n+1}}\) multiplica por \(B\) todas
las contribuciones anteriores y añade \(c_{n+1}\) al acumulado.
Se obtiene la recurrencia y, por sustitución, ambas coordenadas de la
etapa siguiente. La suma elimina \(C_n\).
\end{proof}

Cuando los acarreos son canónicos,
\begin{equation}
 0\leq C_n\leq B^n-1.
 \label{lib03:intervalo-acumulado}
\end{equation}
En particular, \(x_0\geq1\), \(y_0\geq0\) garantizan
\(x_n\geq B^n(x_0-1)+1\) e \(y_n\geq0\) en todas las etapas.

\begin{corollary}[Recuperación de la palabra completa]
\label{lib03:cor-palabra}
Conocidos \(B,n\) y la pareja final, una historia canónica finita
recupera su entrada y todos sus acarreos mediante
\begin{align}
 C_n&=y_n\bmod B^n,&
 y_0&=\left\lfloor\frac{y_n}{B^n}\right\rfloor,&
 x_0&=\frac{x_n+C_n}{B^n},
 \label{lib03:recuperacion-inicial}\\
 c_j&=\left\lfloor\frac{C_n}{B^{n-j}}\right\rfloor\bmod B,
 &&1\leq j\leq n.
 \label{lib03:recuperacion-digitos}
\end{align}
\end{corollary}

\begin{proof}
La cota \eqref{lib03:intervalo-acumulado} hace de
\(y_n=B^ny_0+C_n\) una división euclídea. Su cociente y residuo
proporcionan \(y_0,C_n\); la primera coordenada da \(x_0\).
Las divisiones sucesivas de \(C_n\) por potencias de \(B\) recuperan
los dígitos. Equivalentemente se aplica
\eqref{lib03:residuo-inversor} desde el último paso hasta el primero.
\end{proof}

La profundidad conserva los ceros iniciales de la palabra, aunque estos
no alteren el entero \(C_n\). Si se guarda \(M_0\) en lugar de
\(n\), la igualdad exacta \((x_n+y_n)/M_0=B^n\) recupera la
profundidad. Normalizar la pareja para que sume uno elimina ese contador
de escala cuando no se conserva por otra lectura.

\subsection{Composición de bloques y compatibilidad con la partición}
\label{lib03:composicion}

Un bloque de \(n\) operaciones, con acumulado \(C\), actúa por
\begin{equation}
 L_C^{[n]}(x,y)=(B^nx-C,B^ny+C).
 \label{lib03:bloque}
\end{equation}
La composición de dos bloques satisface
\begin{equation}
 L_D^{[m]}\circ L_C^{[n]}=L_{B^mC+D}^{[n+m]}.
 \label{lib03:ley-bloques}
\end{equation}
En efecto, sustituir la primera imagen en la segunda multiplica su
acumulado por \(B^m\) y añade \(D\). La ley
\[
 (n,C),(m,D)\longmapsto(n+m,B^mC+D)
\]
es asociativa por composición de aplicaciones. Expresa la contribución
del primer bloque en la escala del segundo y mantiene la salida al
reagrupar los mismos eventos ordenados. Agregar eventos nuevos modifica
el acumulado según esta misma ley.

\subsection{Convergencia y control cuantitativo de la lectura}
\label{lib03:limite}

Escribamos \(M_n=B^nM_0\) y
\begin{equation}
 f_n=\frac{x_n}{M_n},\quad g_n=\frac{y_n}{M_n},\quad
 \theta_n=\frac{C_n}{B^n}=\sum_{j=1}^n\frac{c_j}{B^j}.
 \label{lib03:lecturas-normalizadas}
\end{equation}
Las identidades exactas son
\begin{equation}
 f_n=\frac{x_0}{M_0}-\frac{\theta_n}{M_0},\qquad
 g_n=\frac{y_0}{M_0}+\frac{\theta_n}{M_0},\qquad f_n+g_n=1.
 \label{lib03:unidad-normalizada}
\end{equation}

\begin{theorem}[Convergencia de acarreos acotados]
\label{lib03:teo-limite}
Si \(|c_j|\leq C\) para todo \(j\), la serie
\(\theta=\sum_{j\geq1}c_jB^{-j}\) converge absolutamente. Las
lecturas \(f_n,g_n\) convergen a \(f,g\), con \(f+g=1\), y
\begin{equation}
 |f-f_n|=|g-g_n|\leq\frac{C}{M_0(B-1)B^n}.
 \label{lib03:cota-cola}
\end{equation}
\end{theorem}

\begin{proof}
La cola absoluta satisface
\[
 \sum_{j>n}\frac{|c_j|}{B^j}
 \leq C\sum_{j>n}B^{-j}=\frac{C}{(B-1)B^n}.
\]
La convergencia de la serie geométrica y
\eqref{lib03:unidad-normalizada} prueban todas las afirmaciones.
\end{proof}

Para acarreos de ambos signos, una condición suficiente de positividad
en todas las etapas es \(x_0,y_0\geq C/(B-1)\). Para el alfabeto
canónico basta \(x_0\geq1,y_0\geq0\). En este segundo caso,
\(0\leq\theta\leq1\), \(f_n\) decrece y \(g_n\) crece, y
\begin{equation}
 0\leq f_n-f=g-g_n\leq\frac1{M_0B^n}.
 \label{lib03:cota-canonica}
\end{equation}
La suma permanece fija y las proporciones individuales evolucionan.
La historia finita sigue siendo recuperable cuando se conservan las
coordenadas enteras y la escala.

\subsection{Cilindros compatibles y multiplicidad de frontera}
\label{lib03:cilindros}

Consideremos el alfabeto canónico completo y una entrada
\(x_0\geq1,y_0\geq0\). Un prefijo de profundidad \(n\), con
acumulado \(C_n\), admite exactamente los valores límite
\begin{equation}
 I_n(C_n)=\left[\frac{C_n}{B^n},\frac{C_n+1}{B^n}\right]
 \label{lib03:cilindro}
\end{equation}
de la lectura \(\theta\).

\begin{theorem}[Cilindros anidados y distinción de historias]
\label{lib03:teo-cilindros}
Los \(B\) cilindros sucesores de \(I_n(C_n)\) cubren el predecesor y se
solapan únicamente en extremos. Los cilindros de una historia infinita
tienen diámetro \(B^{-n}\) y una intersección de un solo punto.
Dos historias canónicas diferentes producen el mismo valor si y sólo
si, en su primera posición distinta, sus dígitos difieren en uno y las
colas posteriores son, respectivamente, todo \(B-1\) para el dígito
menor y todo cero para el mayor.
\end{theorem}

\begin{proof}
El sucesor asociado a \(c\) tiene extremos
\((BC_n+c)/B^{n+1}\) y \((BC_n+c+1)/B^{n+1}\).
Estas expresiones demuestran inclusión, cobertura y solapamiento sólo
en extremos. La cola de una expansión recorre el intervalo
\([0,B^{-n}]\); la elección sucesiva de uno de los sucesores produce una
expansión de cada punto de ese intervalo. Los intervalos cerrados
anidados tienen diámetros que tienden a cero, lo que prueba existencia
y unicidad de la intersección para cada historia.

Sean \(k\) la primera posición distinta y \(d_k>c_k\). Si las
lecturas coinciden,
\[
 \frac{d_k-c_k}{B^k}
 =\sum_{j>k}\frac{c_j-d_j}{B^j}
 \leq\sum_{j>k}\frac{B-1}{B^j}=\frac1{B^k}.
\]
La diferencia positiva entera \(d_k-c_k\) debe ser uno. La igualdad
en la cota de la cola exige \(c_j=B-1\), \(d_j=0\) para todo
\(j>k\). A la inversa, esas colas tienen diferencia exactamente
\(B^{-k}\), que compensa el primer dígito. Queda probada la
caracterización completa.
\end{proof}

La imagen del cilindro para \(f\) tiene diámetro
\(1/(M_0B^n)\), coincidente con \eqref{lib03:cota-canonica}.
Las historias \((1,0,0,\ldots)\) y
\((0,B-1,B-1,\ldots)\) producen ambas \(\theta=1/B\), conservando
prefijos y estados enteros finitos distintos. Una marca de frontera
mantiene ambas procedencias. Una convención que excluya colas finalmente
constantes \(B-1\) selecciona una expansión para \(\theta<1\);
el extremo \(\theta=1\) se conserva por separado, pues en el
alfabeto fraccionario corresponde a la cola constante \(B-1\).

Historia de operaciones, sistema de cilindros y valor límite son
objetos relacionados mediante estas aplicaciones. La unicidad de la
lectura de una historia coexiste con la multiplicidad de historias en
los valores de frontera.

\subsection{Acarreos firmados y especificación prospectiva}
\label{lib03:firmados}

Permitir acarreos fuera del alfabeto canónico exige conservar sus
elevaciones. Ya en dos pasos, las palabras
\[
 (c_1,c_2)=(1,-1),\qquad(d_1,d_2)=(0,B-1)
\]
producen \(C_2=B-1\) y la misma pareja final. Ambas son acotadas por
\(B-1\) y pueden preservar la positividad, por ejemplo desde
\((73,27)\) en base diez. La palabra firmada, o las elevaciones de
\eqref{lib03:elevacion}, distingue esa información. El teorema de
convergencia permanece válido para ambas palabras; la inversa canónica
tiene, en cambio, el dominio indicado en su enunciado.

La transición \eqref{lib03:caso-729} fija \(c_1=1\). Una regla
adicional que imponga \(c_j=1\) para todo \(j\) produce
\begin{equation}
 \theta_n=\frac{1-B^{-n}}{B-1},\qquad
 \theta=\frac1{B-1}.
 \label{lib03:acarreo-constante}
\end{equation}
En base diez, desde \((73,27)\), la secuencia es
\[
 (73,27)\longmapsto(729,271)\longmapsto(7289,2711)
 \longmapsto(72889,27111),
\]
y la primera fracción converge a
\begin{equation}
 \frac{73}{100}-\frac1{900}=\frac{164}{225}
 =0{,}728888\ldots.
 \label{lib03:limite-constante}
\end{equation}
La conclusión pertenece a esa regla adicional de acarreos constantes.
Con sólo \(c_1=1\) y prolongaciones canónicas no negativas, el cilindro
del primer paso proporciona exactamente
\begin{equation}
 \frac{b^3-1}{B^3}\leq f\leq\frac{b^3}{B^3}.
 \label{lib03:primer-cilindro}
\end{equation}
En base diez es \([0{,}728,0{,}729]\). Cambiar el signo o el alfabeto
cambia esta colección y requiere conservar la memoria adicional
correspondiente. La presencia de \(729\) en la completación determina
la transición aritmética escrita; una identificación con otro lector de
constantes utiliza además su regla efectiva de generación y su
entrelazamiento.

\subsection{Contenido operativo y transición a las amplitudes}
\label{lib03:operativo}

El procedimiento resultante comprende actualización mediante
\eqref{lib03:lc}, comprobación local por residuo, reconstrucción de
toda palabra canónica mediante \eqref{lib03:recuperacion-digitos},
control de la lectura límite mediante \eqref{lib03:cota-cola} y
reagrupación de eventos mediante \eqref{lib03:ley-bloques}. Cada
operación mantiene el orden y el contador de escala que su inversa
utiliza.

La conservación de suma y la conservación cuadrática poseen formas
diferentes. Para una transferencia normalizada
\(\delta=c/(BM)\),
\begin{equation}
 (f-\delta)^2+(g+\delta)^2-(f^2+g^2)
 =2\delta(g-f)+2\delta^2.
 \label{lib03:variacion-cuadratica}
\end{equation}
Esta expresión determina la variación cuadrática y generalmente es
distinta de cero. El enlace con una isometría se obtendrá sobre las
amplitudes de las etiquetas admisibles, mediante la biyección por
residuos de la \cref{lib05:levantamiento}. Así, el transporte aritmético
y su realización unitaria se componen conservando sus respectivas
cantidades, en lugar de identificarlas por una coincidencia escalar.
